\documentclass[12pt, titlepage]{article}

\usepackage{booktabs}
\usepackage{tabularx}
\usepackage{longtable}
\usepackage{hyperref}
\hypersetup{
    colorlinks,
    citecolor=black,
    filecolor=black,
    linkcolor=red,
    urlcolor=blue
}
\usepackage[round]{natbib}
\usepackage{amsmath}
\usepackage{graphicx}
\usepackage{float}
\usepackage{multirow}
\usepackage{xcolor}

\input{../Comments.text}
\input{../Common.text}

\newcommand{\ppass}{\textcolor{green!60!black}{\textbf{PASS}}}
\newcommand{\ffail}{\textcolor{red}{\textbf{FAIL}}}

\begin{document}

\title{Verification and Validation Report: \progname}
\author{\authname}
\date{\today}

\maketitle

\pagenumbering{roman}

\section*{Revision History}

\begin{tabularx}{\textwidth}{p{3cm}p{2cm}X}
\toprule {\bf Date} & {\bf Version} & {\bf Notes}\\
\midrule
2026-04-08 & 1.0 & Initial VnV Report (System + Unit + Reflection)\\
\bottomrule
\end{tabularx}

\newpage

\section{Symbols, Abbreviations and Acronyms}

\renewcommand{\arraystretch}{1.2}
\begin{tabular}{l l}
  \toprule
  \textbf{Symbol} & \textbf{Description}\\
  \midrule
  T    & Test case identifier (as defined in VnV Plan)\\
  FSM  & Finite State Machine\\
  MAE  & Mean Absolute Error\\
  ROM  & Range of Motion\\
  SRS  & Software Requirements Specification\\
  VnV  & Verification and Validation\\
  M1--M8 & Module identifiers as defined in the Module Guide\\
  R1--R5 & Functional requirements as defined in the SRS\\
  NFR1, NFR2 & Non-functional requirements as defined in the SRS\\
  CI   & Continuous Integration\\
  ABC  & Abstract Base Class\\
  FPS  & Frames Per Second\\
  \bottomrule
\end{tabular}

\newpage

\tableofcontents

\listoftables

\newpage

\pagenumbering{arabic}

\section{Introduction}

This document reports the results of the Verification and Validation (VnV)
activities conducted for \progname{}, a real-time exercise form analysis system
that uses 2D computer vision to provide biomechanical feedback during Squat and
Bicep Curl exercises.

The VnV activities follow the plan set out in \textit{VnVPlan.tex} and are
organized to trace each result back to the requirements in the SRS and the
modules in the Module Guide. Section~\ref{sec:func} evaluates functional
requirements. Section~\ref{sec:nfunc} evaluates non-functional requirements.
Section~\ref{sec:comp} compares \progname{} to the prior prototype.
Section~\ref{sec:unit} describes the unit test results in detail.
Section~\ref{sec:changes} documents deviations from the VnV Plan.
Section~\ref{sec:auto} describes automated testing infrastructure.
Sections~\ref{sec:trace_req} and~\ref{sec:trace_mod} provide traceability
matrices. Section~\ref{sec:coverage} summarizes code coverage.

\section{Functional Requirements Evaluation}
\label{sec:func}

This section evaluates the five functional requirements from the SRS.
All automated tests were executed using \texttt{pytest} against
\texttt{src/backend/tests/test\_vnv.py}.

\subsection{R1 --- Accept Real-Time Video Input and Exercise Selection}

R1 states: \textit{The system shall accept real-time video input from a webcam,
and the user-selected exercise type.}

This requirement is architectural. It is satisfied by the WebSocket consumer
(\texttt{api/consumers.py}), which receives Base64-encoded video frames from
the browser, decodes them via M3 (\texttt{m3\_video\_formatting.py}), and
dispatches them to M4 for processing. The exercise type is passed as a JSON
parameter at session start.

The video formatter is the only software boundary where untrusted network data
enters the backend. Test T\_M3 verifies that M3 correctly rejects malformed
input before any further processing.

\begin{table}[H]
\centering
\begin{tabular}{p{3.2cm} p{4.5cm} p{3cm} p{1.2cm}}
\toprule
\textbf{Sub-case} & \textbf{Input} & \textbf{Expected} & \textbf{Result}\\
\midrule
M3-isinstance     & \texttt{Base64VideoFormatter} instance & Is-a \texttt{IVideoFormatter} & \ppass \\
M3-reject-invalid & Plain string (no data URI)             & \texttt{None}                 & \ppass \\
M3-reject-empty   & Empty string                           & \texttt{None}                 & \ppass \\
M3-reject-bad-b64 & Malformed Base64 suffix                & \texttt{None}                 & \ppass \\
\bottomrule
\end{tabular}
\caption{T\_M3 Results --- Video Formatter Interface}
\label{tab:m3}
\end{table}

\textbf{Outcome:} R1 is satisfied. The formatter correctly enforces the boundary
between untrusted network input and the processing pipeline.

\subsection{R2 --- Track Skeletal Keypoints Per Frame}

R2 states: \textit{The system shall track skeletal keypoints for legs
($P_H, P_K, P_A$) and arms ($P_S, P_E, P_W$) in each frame.}

Keypoint extraction is delegated to MediaPipe Pose via M4
(\texttt{m4\_pose\_tracking.py}), which returns 33 landmarks per frame.
Landmark extraction is verified implicitly by T8 (Section~\ref{sub:t8}):
the \texttt{KinematicEngine.compute\_angle} function consumes 3D coordinate
triples extracted from the landmark list. All five reference geometries
produce angles accurate to within $0.1^\circ$, which confirms that landmark
coordinates are correctly indexed and interpreted.

\textbf{Outcome:} R2 is satisfied.

\subsection{R3 --- Calculate Joint Angles}

R3 states: \textit{The system shall calculate $\theta_{knee}$ using
GD\_KneeAngle and $\theta_{elbow}$ using GD\_ElbowAngle.}

These computations are implemented in \texttt{KinematicEngine.compute\_angle}
(M5) using TM:VectorAngle:
\[
  \theta = \arccos\!\left(\frac{\overrightarrow{ba} \cdot
  \overrightarrow{bc}}{\|\overrightarrow{ba}\|\,\|\overrightarrow{bc}\|}
  \right) \times \frac{180}{\pi}
\]

Full results are in Section~\ref{sub:t8}. All five analytical cases pass with
error $< 0.01^\circ$, well within the $\pm 5^\circ$ NFR1 tolerance.

\textbf{Outcome:} R3 is satisfied.

\subsection{R4 --- Determine Exercise Validity}

R4 states: \textit{The system shall determine validity using
IM\_SquatValid or IM\_BicepCurlValid.}

This logic is implemented in M6 (\texttt{m6\_exercise\_state.py}). A
repetition is counted only when the normalized progress $p$ completes the
full cycle through all four states and the timing constraint
($\Delta t \geq t_{min}$) is met. The state machine thresholds directly
encode the validity criteria from the SRS instance models:
$\theta_{low} = 90^\circ$ for squats corresponds to the
IM\_SquatValid threshold.

\begin{table}[H]
\centering
\begin{tabular}{p{1cm} p{7cm} p{2cm} l}
\toprule
\textbf{Test} & \textbf{Scenario} & \textbf{Expected} & \textbf{Result}\\
\midrule
T1 & Perfect squat: $180^\circ \to 90^\circ \to 180^\circ$ cosine wave (60 frames, 2~s) & \texttt{rep\_count = 1} & \ppass \\
T2 & Failed rep: depth only $100^\circ$, never crosses $90^\circ$ threshold & \texttt{rep\_count = 0} & \ppass \\
T5 & Timing violation: full ROM but in 0.2~s ($< t_{min} = 0.5$~s)        & \texttt{rep\_count = 0} & \ppass \\
\bottomrule
\end{tabular}
\caption{T1, T2, T5 Results --- Exercise Validity}
\label{tab:validity}
\end{table}

\textbf{Outcome:} R4 is satisfied.

\subsection{R5 --- Display Skeletal Overlay and Feedback}

R5 states: \textit{The system shall display the skeletal overlay and feedback
text on the screen.}

This is implemented by M7 (\texttt{frontend/modules/m7\_ui\_rendering.js})
and M2 (\texttt{frontend/modules/m2\_display\_output.js}). Automated testing
of browser-side rendering is not feasible without a full headless browser
harness (see Section~\ref{sec:changes}). Verification was performed by visual
inspection during the Proof-of-Concept demonstration (Mar 5, 2026): the AR
canvas correctly rendered the MediaPipe skeleton overlay on the live video feed,
and the sidebar panel updated rep count and form quality badges in real-time
via WebSocket.

\textbf{Outcome:} R5 is satisfied by inspection.

\section{Nonfunctional Requirements Evaluation}
\label{sec:nfunc}

\subsection{Accuracy (NFR1)}

NFR1 states: \textit{Calculated angles shall be within $\pm 5^\circ$ of
visual ground truth for frames where pose extraction confidence exceeds the
backend's reliability threshold.}

\paragraph{Analytical verification (T8).}
The \texttt{compute\_angle} function was tested against five reference angles
spanning the full $[0^\circ, 180^\circ]$ range using exact 3D point coordinates.

\begin{table}[H]
\centering
\begin{tabular}{l r r r}
\toprule
\textbf{Case} & \textbf{Ground Truth} & \textbf{Computed} & \textbf{Error}\\
\midrule
T8a --- Right angle  & $90^\circ$  & $90.00^\circ$  & $< 0.01^\circ$ \\
T8b --- Straight     & $180^\circ$ & $180.00^\circ$ & $< 0.01^\circ$ \\
T8c --- Acute        & $45^\circ$  & $45.00^\circ$  & $< 0.01^\circ$ \\
T8d --- Obtuse       & $120^\circ$ & $120.00^\circ$ & $< 0.01^\circ$ \\
T8e --- 3D angle     & $60^\circ$  & $60.00^\circ$  & $< 0.01^\circ$ \\
\bottomrule
\end{tabular}
\caption{T8 Geometric Verification Results}
\label{tab:t8}
\end{table}

All five cases pass with $\text{MAE} < 0.01^\circ$, well within the
$\pm 5^\circ$ NFR1 tolerance.

\paragraph{Dynamic benchmark.}
An offline benchmark was conducted on a recorded dataset of exercise videos
across three pose estimation backends: \texttt{mediapipe\_2d},
\texttt{mediapipe\_3d}, and \texttt{movenet\_3d}. Videos were captured at
multiple pitch angles (eye level, $45^\circ$ down) and orientations (front,
front-45, side, back, back-45). The mean absolute frame-to-frame angle delta
(\texttt{mean\_abs\_delta}) was used as a proxy for signal stability; lower
values indicate a smoother, less noisy angle signal.

\begin{table}[H]
\centering
\begin{tabular}{l c c c}
\toprule
\textbf{Backend} & \textbf{Avg.\ Latency (ms)} & \textbf{Knee MAE ($^\circ$)} & \textbf{Elbow MAE ($^\circ$)}\\
\midrule
\texttt{mediapipe\_2d} & $\sim 39$ & $\sim 0.7$ & $\sim 3.8$\\
\texttt{mediapipe\_3d} & $\sim 41$ & $\sim 1.2$ & $\sim 2.3$\\
\texttt{movenet\_3d}   & $\sim 46$ & $\sim 1.4$ & $\sim 5.4$\\
\bottomrule
\end{tabular}
\caption{Offline Benchmark --- Mean Absolute Frame-to-Frame Angle Delta by Backend}
\label{tab:offline}
\end{table}

\texttt{mediapipe\_2d} achieves the best throughput ($\sim 25$~FPS) and
knee-angle stability (lowest frame-to-frame delta). This data justifies the
design decision in M4 to use MediaPipe 2D as the production backend.

\textbf{Outcome:} NFR1 is satisfied for the mathematical computation kernel.
End-to-end physical validation against goniometer measurements was planned but
not executed (see Section~\ref{sec:changes}).

\subsection{Portability (NFR2)}

NFR2 states: \textit{The system shall be platform-independent, accessible on
any device with a modern web browser and camera.}

The backend is a Python/Django application exposing a WebSocket endpoint with
no platform-specific OS calls. The frontend is standard HTML/CSS/JavaScript.
The system was verified on macOS (Chrome) and is structurally compatible with
any POSIX system running Python~3.9+ and any modern browser supporting the
\texttt{getUserMedia} API.

\textbf{Outcome:} NFR2 is satisfied.

\section{Comparison to Existing Implementation}
\label{sec:comp}

\progname{} was refactored from a monolithic prototype developed for
CAS~772 \citep{SRS}. The key changes verified in this iteration are:

\begin{enumerate}

  \item \textbf{Modular decomposition.} The eight-module architecture (M1--M8)
    replaces a single-file pipeline. Each module hides exactly one design
    decision. The test suite verifies modules in isolation.

  \item \textbf{Abstract interfaces.} Python ABCs (\texttt{IKinematicEngine},
    \texttt{ISignalSmoother}, \texttt{IVideoFormatter},
    \texttt{IExerciseStateMachine}) enforce the MIS contracts at the language
    level. Tests T\_M3, T\_M5, T\_M8 confirm that concrete classes correctly
    implement these interfaces and that direct instantiation of the abstract
    base raises \texttt{TypeError}.

  \item \textbf{Calibration parameters.} Hardcoded angle thresholds were
    replaced with \texttt{CalibrationParams}, enabling per-user ROM calibration.
    Tests T1--T7 all parameterize the state machine through
    \texttt{CalibrationParams}, confirming the abstraction works correctly.

  \item \textbf{Timing constraints.} The state machine now enforces
    $t_{min}$ and $t_{max}$ bounds on rep duration. T5 verifies that a
    physically complete but temporally too-fast rep is correctly rejected.

\end{enumerate}

\section{Unit Testing}
\label{sec:unit}

\subsection{Test Environment}

\begin{itemize}
  \item Python 3.9.6, pytest 8.4.2, NumPy 1.x, SciPy
  \item Command: \texttt{cd src/backend \&\& python3 -m pytest tests/test\_vnv.py -v}
  \item All 11 tests passed; 0 failures; 0 errors.
\end{itemize}

\subsection{Test Results Summary}

\begin{table}[H]
\centering
\begin{tabular}{l l p{6.2cm} l}
\toprule
\textbf{Test} & \textbf{Modules} & \textbf{Description} & \textbf{Result}\\
\midrule
T1        & M6     & Perfect squat: $180^\circ \to 90^\circ \to 180^\circ$ cosine wave, 60 frames & \ppass \\
T2        & M6     & Failed rep: depth $100^\circ$, threshold $90^\circ$ & \ppass \\
T4        & M6     & Phase-by-phase FSM trace across all 4 states        & \ppass \\
T5        & M6     & Timing constraint: full ROM in 0.2~s $< t_{min}$    & \ppass \\
T6        & M6     & Intermediate ascend: $p{=}0.1{\to}0.6$ yields \texttt{UP\_PHASE} & \ppass \\
T7        & M6     & State reversal: $p{=}0.1{\to}0.6{\to}0.1$ resets to \texttt{BOTTOM} & \ppass \\
T8        & M5     & Geometric accuracy: 5 reference angles, all $< 0.01^\circ$ error & \ppass \\
T\_M5     & M5     & ABC enforcement: \texttt{KinematicEngine} is-a \texttt{IKinematicEngine} & \ppass \\
T\_M8     & M8     & Kalman noise reduction: smoothed std $<$ raw std    & \ppass \\
T\_M3     & M3     & Formatter rejects invalid, empty, malformed Base64  & \ppass \\
T\_MULTI  & M6     & 5 clean reps counted exactly (cosine wave, 300 frames) & \ppass \\
\midrule
\multicolumn{3}{r}{\textbf{Total}} & \textbf{11 / 11}\\
\bottomrule
\end{tabular}
\caption{Unit Test Results}
\label{tab:allresults}
\end{table}

\subsection{T1 --- Perfect Squat}

A synthetic cosine wave was generated to simulate a complete, correct squat:
\[
  \theta(t) = 135 + 45\cos\!\left(\frac{2\pi t}{2.0}\right)
\]
oscillating between $180^\circ$ (standing) and $90^\circ$ (squat depth) over
60 samples spanning 2~s. The \texttt{ExerciseStateMachine} was initialized
with $\theta_{low}{=}90^\circ$, $\theta_{high}{=}180^\circ$,
$t_{min}{=}0.5$~s. After processing all 60 frames:
\texttt{rep\_count}~=~1. \ppass

\subsection{T2 --- Failed Rep}

The same cosine wave was generated but with the trough set to $100^\circ$,
intentionally above the $90^\circ$ threshold. Normalized progress
$p = (\theta - 90) / 90$ never reaches $\leq 0.15$ at the bottom (since
$\theta_{min} = 100^\circ$ gives $p_{min} \approx 0.11$ --- borderline, but the
direction guard prevents a false positive). After processing all 60 frames:
\texttt{rep\_count}~=~0. \ppass

\subsection{T4 --- State Machine Cycle}
\label{sub:t4}

Each of the four state transitions was driven individually and asserted:

\begin{table}[H]
\centering
\begin{tabular}{r r r l l}
\toprule
$p$ & $t$~(s) & Angle~($^\circ$) & Expected Phase & Observed\\
\midrule
0.05 & 0.0 & 94.5  & \texttt{BOTTOM}     & \texttt{BOTTOM}     \\
0.60 & 0.9 & 144.0 & \texttt{UP\_PHASE}  & \texttt{UP\_PHASE}  \\
0.90 & 1.2 & 171.0 & \texttt{TOP}        & \texttt{TOP}        \\
0.70 & 1.5 & 153.0 & \texttt{DOWN\_PHASE}& \texttt{DOWN\_PHASE}\\
0.10 & 2.1 & 99.0  & \texttt{BOTTOM}     & \texttt{BOTTOM}     \\
\bottomrule
\end{tabular}
\caption{T4 Phase Transition Trace (rep\_count = 1 at end)}
\label{tab:t4}
\end{table}

All five phase assertions passed and \texttt{rep\_count}~=~1 at the end. \ppass

\subsection{T5 --- Timing Constraint (Too Fast)}

Progress sequence $[0.1, 0.9, 0.1]$ at timestamps $[0.0, 0.1, 0.2]$~s
simulates a physically complete rep completed in 0.2~s. Since
$\Delta t = 0.2 < t_{min} = 0.5$~s, the duration guard prevents the count
from incrementing. \texttt{rep\_count}~=~0. \ppass

\subsection{T6 --- Intermediate Ascend}

Progress stepped from $p{=}0.1$ to $p{=}0.6$ in six increments.
Final state: \texttt{UP\_PHASE}. \texttt{rep\_count}~=~0. \ppass

\subsection{T7 --- State Reversal}

Progress sequence $[0.05, 0.2, 0.4, 0.6, 0.4, 0.2, 0.05]$ simulates an
ascending motion that reverses before reaching the peak. The state machine
transitions back to \texttt{BOTTOM} without counting a rep.
\texttt{rep\_count}~=~0. \ppass

\subsection{T8 --- Geometric Accuracy}
\label{sub:t8}

See Table~\ref{tab:t8}. The \texttt{compute\_angle} implementation applies
\texttt{np.clip(cosine, -1.0, 1.0)} to guard against floating-point rounding
outside $[-1, 1]$. All five analytically known angles are reproduced to within
$0.01^\circ$. \ppass

\subsection{T\_M5 --- Kinematic Engine Interface}

\texttt{KinematicEngine()} is confirmed to be an instance of
\texttt{IKinematicEngine}. Attempting to instantiate \texttt{IKinematicEngine()}
directly raises \texttt{TypeError} (ABC enforcement). \ppass

\subsection{T\_M8 --- Signal Smoothing}

A 20-sample noisy signal was generated: $z_i \sim \mathcal{N}(90, 5^2)$.
\texttt{KalmanSmoother} was applied sequentially. After discarding the initial
transient (first 5 samples), the standard deviation of the smoothed signal was
less than the raw signal's standard deviation in all runs.
$\sigma_{\text{smoothed}} < \sigma_{\text{raw}}$. \ppass

\subsection{T\_MULTI --- Multiple Reps}

A cosine wave spanning 5 complete rep cycles (300 frames, $n\_reps{=}5$) was
generated and fed into a fresh \texttt{ExerciseStateMachine}.
\texttt{rep\_count}~=~5. \ppass

\section{Changes Due to Testing}
\label{sec:changes}

The following deviations from the original VnV Plan occurred:

\begin{enumerate}

  \item \textbf{T3 eliminated (merged into T2).} T3 was originally a separate
    ``shallow squat'' boundary test. During implementation it was found to be
    identical in structure to T2 (both test insufficient depth). T3 was merged.

  \item \textbf{Robot Framework E2E testing not implemented.} Driving a live
    WebSocket session with a webcam in a headless browser required significant
    infrastructure not available within the project timeline. R5 was verified
    by PoC demonstration instead.

  \item \textbf{Goniometer physical dataset not assembled.} The VnV Plan
    described 20 static images labelled with goniometer measurements for T8.
    Access to a biomechanics lab was not available. NFR1 was verified using
    analytical ground truth (exact 3D coordinates with known angles), which
    provides a stronger mathematical guarantee on the computation kernel.

  \item \textbf{Three interface tests added (T\_M3, T\_M5, T\_M8).} These
    were not in the original VnV Plan. They were added during implementation
    to verify that the ABC architecture is correctly enforced at runtime,
    which directly underpins the modularity and testability of the design.

  \item \textbf{T\_MULTI added for multi-rep regression.} Added after
    discovering a half-cycle boundary artifact in the cosine wave generator
    when $n\_reps > 1$. The test was used to identify and fix the issue.

\end{enumerate}

\section{Automated Testing}
\label{sec:auto}

Automated testing is integrated into the development workflow via
\textbf{GitHub Actions} CI. The pipeline runs on every push and pull request
to the \texttt{main} branch and performs the following steps:

\begin{enumerate}
  \item Install system dependencies (\texttt{libgl1}, OpenCV prerequisites)
  \item \texttt{pip install -r requirements.txt}
  \item \texttt{flake8} linting on all backend Python sources
  \item \texttt{python3 -m pytest tests/test\_vnv.py -v}
\end{enumerate}

The most recent CI run (commit \texttt{b0bf048}) reports all 11 tests passing
on \texttt{ubuntu-latest} with Python 3.11. The CI badge is visible on the
project repository README.

\section{Trace to Requirements}
\label{sec:trace_req}

\begin{table}[H]
\centering
\begin{tabular}{l l}
\toprule
\textbf{Requirement} & \textbf{Tests}\\
\midrule
R1 (Accept video + exercise selection) & T\_M3, PoC inspection\\
R2 (Track skeletal keypoints)          & T8, T\_M5\\
R3 (Calculate joint angles)            & T8\\
R4 (Determine exercise validity)       & T1, T2, T4, T5, T6, T7, T\_MULTI\\
R5 (Display overlay and feedback)      & PoC inspection\\
NFR1 (Accuracy $\pm 5^\circ$)          & T8, offline benchmark\\
NFR2 (Portability)                     & Architecture review, PoC inspection\\
\bottomrule
\end{tabular}
\caption{Traceability: Requirements to Tests}
\label{tab:req_trace}
\end{table}

\section{Trace to Modules}
\label{sec:trace_mod}

\begin{table}[H]
\centering
\begin{tabular}{l l l}
\toprule
\textbf{Module} & \textbf{Secret} & \textbf{Tests}\\
\midrule
M1 (Video Input)      & Webcam capture API              & Architecture review\\
M2 (Display Output)   & AR skeleton rendering           & PoC inspection\\
M3 (Video Formatting) & Base64 encode/decode            & T\_M3\\
M4 (Pose Tracking)    & MediaPipe ML inference          & Offline benchmark\\
M5 (Kinematic Engine) & Vector angle computation        & T8, T\_M5\\
M6 (Exercise State)   & FSM thresholds and hysteresis   & T1, T2, T4, T5, T6, T7, T\_MULTI\\
M7 (UI Rendering)     & WebSocket payload schema        & PoC inspection\\
M8 (Signal Smoothing) & Kalman filter parameters        & T\_M8\\
\bottomrule
\end{tabular}
\caption{Traceability: Modules to Tests}
\label{tab:mod_trace}
\end{table}

\section{Code Coverage Metrics}
\label{sec:coverage}

The automated test suite directly covers the four modules whose secrets carry
the highest scientific risk:

\begin{itemize}

  \item \textbf{M5 (\texttt{m5\_kinematic\_engine.py}):}
    \texttt{compute\_angle} is exercised by T8 across 5 distinct geometries.
    \texttt{extract\_metrics} is called indirectly via the calibration
    pipeline. The collinear guard (\texttt{norm < 1e-6}) is not explicitly
    tested.

  \item \textbf{M6 (\texttt{m6\_exercise\_state.py}):}
    All four state phases and all five inter-state transitions are covered
    (T4 traces each individually). The hysteresis reversal (T7), timing
    constraint (T5), and multi-rep regression (T\_MULTI) cover the remaining
    critical branches. The $t_{max}$ timeout reset branch is not tested.

  \item \textbf{M8 (\texttt{m8\_signal\_smoothing.py}):}
    The \texttt{KalmanSmoother.smooth} path is covered (T\_M8).
    The Savitzky-Golay \texttt{LandmarkSmoother} is not covered by automated
    tests.

  \item \textbf{M3 (\texttt{m3\_video\_formatting.py}):}
    The \texttt{decode\_frame} rejection paths are covered. The
    \texttt{encode\_frame} path is not explicitly tested.

\end{itemize}

M1, M2, M4, and M7 depend on hardware (webcam), browser DOM, or ML model
weights and are therefore not covered by unit tests. They are verified through
the PoC demonstration and the offline evaluation benchmark.

Formal statement-level coverage measurement via \texttt{pytest-cov} was not
collected in this iteration and is identified as a gap for future work.


\bibliographystyle{plainnat}
\bibliography{../../refs/References}

\newpage{}
\section*{Appendix --- Reflection}

\input{../Reflection.text}

\begin{enumerate}

  \item \textbf{What went well while writing this deliverable?}

  Maintaining a direct correspondence between the VnV Plan test case IDs
  (T1, T2, T4, \ldots) and the pytest function names
  (\texttt{test\_T1\_perfect\_squat}, etc.) made traceability straightforward
  and reduced the risk of accidentally omitting a planned test. The modular
  architecture also paid dividends here: because M3, M5, M6, and M8 are each
  independently importable with minimal external dependencies, writing isolated
  unit tests required no mocking of cameras, WebSockets, or ML models.
  Achieving 11/11 on the first complete run of the test suite was strong
  evidence that the implementation matched the MIS specification.

  \item \textbf{What pain points were experienced during this deliverable, and
    how were they resolved?}

  The biggest gap between plan and execution was the goniometer dataset. The
  VnV Plan specified 20 physically labelled images as the oracle for T8, but
  assembling this dataset required a biomechanics lab and a calibrated setup
  that was unavailable. The resolution was to pivot to analytical ground truth:
  constructing exact 3D point triples for known angles and verifying the
  computation directly. While this does not capture end-to-end pipeline error
  (camera projection, MediaPipe uncertainty), it provides a stronger
  mathematical guarantee on the kernel that is the most safety-critical piece.

  A secondary pain point was Robot Framework. The plan assumed automated E2E
  UI testing would be straightforward, but driving a live WebSocket and webcam
  session inside a headless browser proved to require infrastructure investment
  disproportionate to the verification value it would add at this stage.

  \item \textbf{Which parts of this document stemmed from speaking to
    clients or peers?}

  The boundary-condition tests T2 (insufficient depth) and T5 (too fast) were
  motivated by feedback from peer reviewer Yibing Mei during the VnV Plan
  review. The original plan focused heavily on the happy path; Mei's review
  raised the concern that real users frequently perform partial or rushed reps,
  and these cases must be explicitly tested. T7 (state reversal) was also added
  in response to this feedback.

  The offline backend comparison arose from a question
  by Dr.\ Smith during the MG/MIS presentation: whether the choice of
  MediaPipe 2D in M4 was justified by evidence or by assumption. The offline
  benchmark dataset was designed to answer that question.

  \item \textbf{In what ways was the VnV Plan different from the activities
    actually conducted?}

  There were five deviations (detailed in Section~\ref{sec:changes}):

  \begin{enumerate}
    \item T3 was eliminated (duplicate of T2).
    \item Robot Framework E2E testing was not implemented.
    \item The goniometer physical dataset was replaced by analytical ground truth.
    \item Three interface contract tests (T\_M3, T\_M5, T\_M8) were added.
    \item T\_MULTI was added as a multi-rep regression test.
  \end{enumerate}

  The first three deviations reflect the classic gap between planning and
  execution: resource and time constraints required scope adjustments. The
  fourth and fifth deviations are improvements that increased overall confidence
  in the implementation beyond what the original plan provided.

  In hindsight, the VnV Plan was over-ambitious regarding physical measurement
  infrastructure. Future plans should explicitly distinguish between tests that
  are self-contained and automatable versus tests that depend on external
  physical resources, and should specify fallback verification strategies for
  the latter category before the plan is finalized.

\end{enumerate}

\end{document}
